\chapter{สถาปัตยกรรมเว็บสมัยใหม่และการประมวลผล 60 FPS}

เว็บเบราว์เซอร์ยุคใหม่มีขีดความสามารถเทียบเท่าเนทีฟแอปพลิเคชัน การนำเทคโนโลยี CSS ยุคใหม่ร่วมกับเอนจินกราฟิกส์ ช่วยให้สร้างอินเทอร์เฟซที่สวยงามและตอบสนองได้แบบทันทีทันใด

\section{เทคโนโลยี CSS ยุคใหม่}

การออกแบบเลย์เอาต์ที่ยืดหยุ่นในปัจจุบัน ไม่จำเป็นต้องพึ่งพา JavaScript ในการดักจับขนาดหน้าจอ:
\begin{itemize}
    \item \textbf{Container Queries (\texttt{@container}):} ปรับเปลี่ยนสไตล์ตามขนาดของการ์ดหรือพาเนล ไม่จำกัดเฉพาะขนาดหน้าต่างเบราว์เซอร์
    \item \textbf{Parent Selector (\texttt{:has()}):} ควบคุมสไตล์ของพาเรนต์ตามสถานะขององค์ประกอบลูก
    \item \textbf{View Transitions API:} สร้างแอนิเมชันการสลับหน้าระหว่างการทดลองที่ลื่นไหลไร้รอยต่อ
\end{itemize}

\section{การป้องกันปัญหา Garbage Collection Trashing}

ในแอนิเมชันความเร็ว 60 เฟรมต่อวินาที แต่ละเฟรมมีเวลาประมวลผลเพียง $16.67\,\text{ms}$:

\begin{equation}
\Delta t_{\text{frame}} = \frac{1000\,\text{ms}}{60\,\text{FPS}} \approx 16.667\,\text{ms}
\end{equation}

หากโค้ดสร้างออบเจกต์ชั่วคราวมากเกินไป ตัวเก็บขยะ (Garbage Collector) จะแทรกการทำงานและทำให้เวลาประมวลผลเฟรมเกิน $16.67\,\text{ms}$ เกิดอาการภาพกระตุก (Jank) การใช้ ArrayBuffer ชนิด TypedArray และ Object Pool ช่วยขจัดปัญหานี้ได้อย่างสิ้นเชิง
